\section{Coeficiente mixto y paso a las cuatro actividades físicas}
Restamos a $R$ su constante y el campo $r\cdot x$, y a $A$ su constante. Escribimos
$R_c,A_c$ para las funciones centradas. Las cuatro actividades auxiliares son $\zeta_i$.
En el punto físico $\zeta_i=e^{r_i}=O(s)$. En efecto, sin enlaces extra, las cuatro
marcas locales son Bernoulli en el campo de la silla y su óptimo son las marginales de
Bethe $P_i$ en la máscara $A$ original. Luego
$\log(P_i/(1-P_i))=r_i+\partial_iR_c(P)$, y $\partial_iR_c(P)=O(s^2)$.
Si un sitio está ausente, la diferencia mixta que lo requiere se anula.
Para el conjunto de enlaces $E$, sea $\chi_E(X)$ el indicador de las exclusiones locales.
La parte exacta normalizada es $\log\sum_X\chi_E(X)\zeta^X e^{g(X)}$.
Para el predictor sea
\[
 \mathcal K(\lambda)=\log\sum_X\chi_E(X)\zeta^X e^{\lambda\cdot X},\quad
 x=\nabla\mathcal K(\lambda),\quad V=\mathcal K''(\lambda),\quad\lambda=\nabla R_c(x).
\]
El valor normalizado, por el apartado 2 y la dualidad del bosque, es
\[
 \mathcal K(\lambda)-\lambda\cdot x+R_c(x)+A_c(x)-\tfrac12\log\det(I-VR''(x)).
\]
En $\zeta$ complejo estos son gérmenes holomorfos de las mismas fórmulas reales.
En un polidisco $|\zeta_i|\le\rho$ pequeño e independiente de $j$, el átomo vacío
tiene peso uno, las otras masas son pequeñas y las derivadas de $\mathcal K$ son acotadas.
La ecuación de $x$ se resuelve en una bola fija por contracción con constante $O(s)$.
Las ramas de logaritmos y del determinante son uniformes, incluido sobre los ejes.
En esta bola sea $x_0=\nabla\mathcal K(0)$ y $V_0=\mathcal K''(0)$. La expansión del
predictor vale, uniformemente en el polidisco,
\[
 \mathcal K(0)+a[x_0]+\tfrac12 H[x_0,x_0]+\tfrac12\tr(V_0H)+O(s^2).
\]
El término lineal de $\mathcal K(\lambda)$ cancela el de $-\lambda\cdot x$ a ese orden.
El mismo desarrollo para el exacto se obtiene del primer resultado sobre $g(X)$
por $V_0+x_0x_0^t=\E[XX^t]$ en la ley local (identidad algebraica para pesos complejos).
En particular el error de cada caso cuesta $O(s^2)$ en un radio fijo.

\paragraph{Cálculo del único coeficiente de grado cuatro.}
En la diferencia mixta $D$ de los errores, poner $\zeta_i=0$ hace inoperante el
enlace correspondiente. Así $D$ es divisible por $\prod_i\zeta_i$. Para su coeficiente
cuadrilineal, en cada componente sustituimos la masa del átomo completo por una variable
independiente, $S$ en el primer par y $T$ en el segundo. La doble diferencia conserva
solamente términos que usan ambas variables. Al volver a $S=\zeta_1\zeta_2$,
$T=\zeta_3\zeta_4$, cualquier otro sitio introduciría un grado mayor.
Basta pues tomar los átomos $0,u,v_*,u+v_*$ de pesos $1,S,T,ST$ en la ley base.
El átomo $ST$ es necesario. Trabajamos módulo $S^2,T^2$; en lo siguiente $v=v_*$.
El coeficiente exacto es
\[
 e^{g(u)+g(v)}(e^{\Delta_{u,v}g}-1).
\]
En el predictor la base factoriza entre los dos grupos. La ecuación de la silla da
$x=Su+Tv+STh(u+v)$ y $V=Suu^t+Tvv^t+STh(uu^t+vv^t)$.
La parte de acción contribuye $h$ al coeficiente. Las derivadas del término $A_c(x)$,
de $\tr(VR'')/2$ y de $\tr[(VR'')^2]/4$ dan en total con la acción
\[
 h+A''(0)[u,v]+h\,a[u+v]
 +\tfrac12\{h(H[u,u]+H[v,v])+U_3[u,u,v]+U_3[u,v,v]+h^2\}.
\]
Las trazas superiores desaparecen pues cada factor $V$ tiene grado positivo.
Como $H,a=O(s)$ y $U_3,A''=O(s^2)$, aplicar los dos resultados del apartado 7
a la expresión exacta da el mismo coeficiente hasta $O(s^3)$.
Por Cauchy en el polidisco fijo, tras dividir $D$ por $\prod_i\zeta_i$, todas las
derivadas primeras son $O(s^2)$ en radio menor, por la cota uniforme anterior.
Su valor en el origen es el coeficiente $O(s^3)$ que acabamos de comprobar.
En las actividades físicas $O(s)$ se obtiene $|D|\le C s^7$.
Las constantes que se retiraron son comunes a los cuatro casos; por tanto $D=w_{bd}$
en el punto físico. No se divide por una marginal física potencialmente pequeña.
